% =============================================================================
% Section 4: Event Selection
% =============================================================================

\section{Event Selection}
\label{sec:selection}

% AGENT: Fill from signal lead Wave 2 report
% Document the event selection strategy, optimization, and categorization.

\subsection{Preselection}
\label{sec:selection:preselection}

% AGENT: Fill preselection criteria from SelectionEngine configuration
% (gad/config/selection/*.yaml). Include trigger matching, object multiplicity,
% and basic kinematic requirements.

A common preselection is applied to all events to ensure basic analysis
requirements are met. The preselection criteria are:
\begin{itemize}
  \item Trigger: [trigger requirement]
  \item Object multiplicity: [requirements]
  \item Basic kinematics: [requirements]
\end{itemize}

\subsection{Optimized Selection}
\label{sec:selection:optimized}

% AGENT: Fill optimized selection cuts from Wave 2 optimization results.
% Include the optimization metric used (e.g., Asimov significance).

The selection is further optimized to maximize the expected signal sensitivity.
The optimization is performed using [metric] as the figure of merit.
Table~\ref{tab:cutflow} summarizes the event yields after each selection
requirement.

% AGENT: Fill cut-flow table from SelectionEngine.apply_cuts() output.
% Include data (if sideband), signal MC, and background MC yields.

\begin{table}[htbp]
  \centering
  \caption{Cut-flow table showing event yields after each selection requirement.
           Signal yields are normalized to [normalization].}
  \label{tab:cutflow}
  \begin{tabular}{lrrrrr}
    \toprule
    Selection & Signal & Bkg 1 & Bkg 2 & Bkg 3 & Total Bkg \\
    \midrule
    % AGENT: Fill rows from cut-flow results
    Preselection       & -- & -- & -- & -- & -- \\
    Cut 1              & -- & -- & -- & -- & -- \\
    Cut 2              & -- & -- & -- & -- & -- \\
    Final selection    & -- & -- & -- & -- & -- \\
    \bottomrule
  \end{tabular}
\end{table}

\subsection{BDT Discriminant}
\label{sec:selection:bdt}

% AGENT: Fill BDT configuration and performance from MVA Wave 2 results.
% Include input variables, training configuration, ROC curve, and
% feature importance ranking.

A Boosted Decision Tree (BDT) classifier is trained to separate signal from
background. The BDT uses [N] input variables, selected based on their
separation power and stability. The input variables and their rankings are
listed in Table~\ref{tab:bdt_variables}.

% AGENT: Fill BDT variable table from BDTTrainer feature importance results.

\begin{table}[htbp]
  \centering
  \caption{BDT input variables ranked by importance (gain).}
  \label{tab:bdt_variables}
  \begin{tabular}{llr}
    \toprule
    Rank & Variable & Importance \\
    \midrule
    % AGENT: Fill from BDTTrainer.feature_importance output
    1 & [variable] & [importance] \\
    2 & [variable] & [importance] \\
    3 & [variable] & [importance] \\
    \bottomrule
  \end{tabular}
\end{table}

% AGENT: Include N-1 plot figure placeholders.
% Reference the plot files generated by SelectionEngine.plot_n_minus_1().

\begin{figure}[htbp]
  \centering
  % AGENT: Replace with actual N-1 plot paths from Wave 2 output
  % \includegraphics[width=0.45\textwidth]{figures/n_minus_1_var1.pdf}
  % \includegraphics[width=0.45\textwidth]{figures/n_minus_1_var2.pdf}
  \caption{N-1 distributions for key selection variables. Data shown in
           sideband regions only (signal region is blinded).}
  \label{fig:n_minus_1}
\end{figure}

\subsection{Event Categorization}
\label{sec:selection:categorization}

% AGENT: Fill categorization scheme if applicable.
% Document how events are split into analysis categories/channels.

Events passing the final selection are categorized into [N] categories
based on [categorization criteria], designed to exploit differences in
signal-to-background ratio across the phase space.
